Medical image segmentation, as a core task in medical image analysis, plays an irreplaceable role in clinical procedures such as disease diagnosis, surgical planning, lesion assessment, and treatment evaluation. Howevercurrent mainstream deep learning-based medical image segmentation methods typically rely on large amounts of high-quality, finely annotated training data. In practical clinical environments, acquiring such annotated data faces numerous challenges: the annotation process is time-consuming and labor-intensive, requires specialized medical knowledge, and there are significant data distribution differences across medical institutions. These issues result in high annotation costs and limited data scale. Furthermore, existing models often perform well on specific datasets but exhibit poor generalization capabilities across institutions, modalities, and anatomical sites, severely restricting their promotion and application in real-world medical scenarios.

Therefore, this study aims to address the practical bottleneck of annotation scarcity by systematically constructing a theoretical and methodological framework for adaptive medical image segmentation with strong generalization capabilities. Specifically, this study first constructs a cognitive uncertainty-guided semi-supervised learning framework based on the CA dataset, aiming to improve model performance by quantifying the uncertainty of predictions within the Mean Teacher architecture to enhance pseudo-label quality. Subsequently, regarding network architecture design, this research deeply investigates different fusion paradigms of U-Net and Transformers, and further introduces a channel attention mechanism into the skip connections to strengthen the model's ability for adaptive selection and integration of multi-scale features.
